% =============================================================================
% 03_2_hauptteil_cloud_config.tex — Domäne: Cloud-Konfiguration
% =============================================================================
% Dieser Abschnitt dokumentiert die Cloud-Infrastruktur und Backend-
% Konfiguration des Geotracker-Projekts. Gliederung gemäss Schulvorgabe:
%   1. Analyse und Design
%   2. Softwaredokumentation
%   3. Reflexion
% =============================================================================

\section{Cloud-Konfiguration}
\label{sec:cloud}

% Kurze Einführung in die Cloud-Domäne
% Welche Cloud-Plattform wird genutzt? Welche Dienste sind im Einsatz?
% (z. B. MQTT-Broker, Datenbank, API-Server)
<<Kurze Einführung in die Cloud-Architektur (z.\,B.\ AWS IoT Core, Azure IoT Hub,
selbst gehosteter MQTT-Broker, Datenbank) und deren Rolle im Gesamtsystem.>>

% --- 1. Analyse und Design ---------------------------------------------------
\subsection{Analyse und Design}

\subsubsection{Anforderungsanalyse}
% Welche Anforderungen hat das Cloud-Backend zu erfüllen?
% (Skalierbarkeit, Verfügbarkeit, Datenschutz, Kosten)
<<Anforderungen an das Cloud-Backend auflisten (z.\,B.\ Empfang von GPS-Daten
per MQTT, Persistierung in Datenbank, REST-API für das Web-Frontend,
Zugriffskontrolle).>>

\subsubsection{Architekturdesign und Diagramme}
% Skizziere die Cloud-Architektur (Komponenten, Datenfluss).
% Füge ein Architektur- oder Sequenzdiagramm ein.
% \begin{figure}[htbp]
%   \centering
%   \includegraphics[width=0.85\textwidth]{uml_cloud_architektur}
%   \caption{Architekturdiagramm der Cloud-Konfiguration}
%   \label{fig:cloud_architektur}
% \end{figure}
<<Architektur- und Sequenzdiagramme hier einfügen und den Datenfluss von der
Firmware bis zur Web-Anwendung erläutern.>>

% --- 2. Softwaredokumentation ------------------------------------------------
\subsection{Softwaredokumentation}

\subsubsection{Beschreibung und Konzepte}
% Dokumentiere die eingesetzten Dienste, Protokolle und Konfigurationen.
% Wie ist der MQTT-Broker konfiguriert? Welche Datenbankstruktur wird verwendet?
<<Dienste und Konfigurationen beschreiben (z.\,B.\ MQTT-Topics, Datenbankschema,
API-Endpunkte, Authentifizierungsmechanismus). Relevante Konfigurationsdateien
oder Codeausschnitte als Listings einfügen.>>

% Beispiel:
% \begin{lstlisting}[caption={MQTT-Topic-Struktur}, label={lst:mqtt_topics},
%                    language={}]
%   geotracker/<device-id>/location   % GPS-Koordinaten
%   geotracker/<device-id>/status     % Gerätestatus
% \end{lstlisting}

\subsubsection{Schwierigkeiten und Lösungsansätze}
% Welche Herausforderungen traten bei der Cloud-Konfiguration auf?
% (z. B. Latenz, Authentifizierung, Skalierung, Kosten)
<<Technische Herausforderungen dokumentieren (z.\,B.\ TLS-Zertifikatsverwaltung,
Datenbank-Performance bei vielen GPS-Punkten) und Lösungsstrategien beschreiben.>>

\subsubsection{Testkonzept}
% Wie wurde das Cloud-Backend getestet?
% (End-to-End-Tests, Lasttests, Sicherheitstests)
<<Testkonzept für das Cloud-Backend beschreiben. Testergebnisse und
verwendete Tools dokumentieren (z.\,B.\ MQTT-Explorer, Postman für API-Tests).>>

% --- 3. Reflexion ------------------------------------------------------------
\subsection{Reflexion}
% Bewerte die Cloud-Lösung kritisch.
% Würdest du dieselbe Plattform wieder wählen? Was würdest du verbessern?
<<Reflexion zur Cloud-Konfiguration: Stärken (z.\,B.\ gute Skalierbarkeit),
Schwächen (z.\,B.\ hohe Komplexität, Kosten) und Verbesserungsvorschläge.>>
